\section{간단한 게임}

St. Petersburg Game과 같은 의미를 지니지만 유한한 기대값을 가지는 게임을 통해 기대값 기반의 이론과 상식이 어긋나는 경우를 보여줄 수 있다. 참가자는 1/4 확률로 40원을 얻거나 3/4 확률로 0원을 얻는다. 이때 참가자가 얻는 상금의 기대값은 10원이다.

참가자는 참가 비용으로 상금의 기대값보다 낮은 금액을 지불한다면 게임에 참여하는 것이 합리적이다. 기대값 기반의 합리적인 이유를 다른 각도에서 바라보면, 참가자가 게임을 반복해서 참여했을 때 얻는 총 상금이 총 참가 비용보다 클 확률이 높기 때문이다. 즉 참가자가 본전을 찾을 확률이 높다면 게임에 참여하는 것이 합리적이다.

\includegraphics[width = \linewidth]{simple_game/E=2.png}

위 그림은 참가자가 게임을 반복해서 참여했을 때 본전을 찾을 확률을 나타낸 것이다. 참가자가 상금의 기대값보다 작은 금액을 지불한다면, 게임을 1000번 이상 참여했을 때 1에 가까운 확률로 돈을 잃지 않는다는 것을 알 수 있다. 참가자가 상금의 기대값과 같은 금액을 지불한다면, 본전을 찾을 확률이 50\%로 수렴하는 것을 볼 수 있다.


\section{희박한 확률 아래에서의 간단한 게임}

이제 간단한 게임을 변형하여 참가자가 $1/2^{10}$ 확률로 $10 \times 2^{10}$원을 얻거나 0원을 얻는 게임을 살펴보자. 이때 참가자가 얻는 상금의 기대값은 여전히 10원이다.

\includegraphics[width = \linewidth]{simple_game/E=10.png}

위 그림에서 참가자가 본전을 되찾을 확률은 상금의 기대값과 참가 비용에 따라서 여전히 1, 0.5, 0으로 수렴하려는 경향을 확인할 수 있다. 그러나 성공 확률이 1/4 인 게임보다 본전을 되찾을 확률이 수렴하는 속도가 느려진 것을 확인할 수 있다. 특히 게임을 100회 연속으로 참여하는 경우, 참가자가 1회도 성공하지 못하고 모든 참가 비용을 소모할 확률은 참가 비용과 상관 없는 것을 확인할 수 있다.


\includegraphics[width = \linewidth]{simple_game/cost=10.png}
위 그림은 상금의 기대값과 참가 비용이 같음에도 불구하고, 성공 확률이 희박해질 수록 참가자가 본전을 되찾을 확률이 0.5로 수렴하는 속도가 느려지는 것을 보인다.
